\chapter{Conclusiones y Trabajo Futuro}
\label{chap:conclusiones}

\section{Conclusiones Generales}
\label{sec:conclusiones_generales}

Este Trabajo de Fin de Máster ha abordado el problema de la selección de carteras financieras bajo restricciones
discretas de cardinalidad mediante un algoritmo cuántico híbrido basado en QAOA, comparando su comportamiento
frente a un solucionador exacto clásico (Gurobi) y frente a una metaheurística de referencia industrial (Recocido
Simulado). Más allá de la implementación técnica, el propósito último del proyecto era responder a una pregunta
concreta: \emph{¿puede un ansatz variacional cuántico, diseñado específicamente para respetar restricciones
combinatorias duras, ofrecer un compromiso competitivo entre viabilidad, precisión y robustez financiera frente a
los métodos clásicos establecidos?}

Los resultados recogidos en el Capítulo 6 permiten responder afirmativamente a esta pregunta, con matices
importantes. La sustitución de las penalizaciones QUBO blandas por el confinamiento geométrico del \emph{XY-Mixer}
sobre el estado de Dicke ha resuelto de forma estructural el problema de viabilidad: mientras que el QAOA estándar
degrada progresivamente su tasa de factibilidad a medida que crece el universo de activos $N$, la arquitectura
XY-QAOA mantiene una viabilidad matemática del 100\,\% en todas las instancias evaluadas, por construcción y no
por ajuste estadístico. En segundo lugar, la combinación de inicialización adiabática guiada (TQA)ha permitido mantener el Optimization GAP por debajo del umbral objetivo del 5\,\% en la práctica totalidad
del rango de escalado estudiado ($N=6$ a $N=20$), con la única excepción de un repunte leve y explicado (hasta
$\sim\!4,8\,\%$) en las instancias más densas, atribuible al límite de expresividad de un circuito de profundidad
fija. Por último, la evaluación financiera fuera de muestra ha revelado un hallazgo especialmente relevante desde
el punto de vista práctico: el solucionador exacto clásico, pese a optimizar perfectamente en la muestra de
entrenamiento, sufre una caída de Ratio de Sharpe por sobreajuste al enfrentarse a datos de mercado no vistos,
mientras que la solución cuántica regularizada generaliza mejor, ofreciendo carteras más robustas en condiciones
de mercado reales.

En conjunto, estas evidencias respaldan la tesis central del proyecto: en el régimen NISQ actual, un algoritmo
cuántico variacional bien diseñado no aventaja a los métodos clásicos en velocidad ni en optimalidad matemática
estricta, pero sí puede aportar valor diferencial en viabilidad estructural y en robustez financiera fuera de
muestra, dos propiedades tan relevantes para la práctica del sector Fintech como la propia optimalidad.

\section{Grado de Cumplimiento de los Objetivos}
\label{sec:cumplimiento_objetivos}

Contrastando los resultados del Capítulo \ref{ch:resultados} con los objetivos puestos en el capitulo \ref{sec:objetivos_especificos}, los cuatro objetivos
específicos se consideran cumplidos. OE1 queda validado mediante la equivalencia analítica estricta entre el
modelo QUBO y el Hamiltoniano de Ising (sección \ref{sec:verificacion}). OE2 se cumple con holgura: el XY-QAOA alcanza una
viabilidad del 100\,\% en todas las instancias evaluadas ($N=6$ a $N=20$), frente a la degradación progresiva
del QAOA estándar (Figura 6.1). OE3 también se cumple en términos del criterio cuantitativo fijado ($\geq 15\,\%$
de reducción de variabilidad del GAP), con una reducción global de \textbf{76.3\%}; cabe matizar,
no obstante, que la ablación de la Sección 6.1 muestra que esta mejora procede principalmente de la
inicialización TQA y no de la penalización Ridge L2 aislada, un matiz que se documenta por rigor científico y se
retoma en la Sección \ref{sec:limitaciones}. Por último, OE4 se cumple: el GAP medio se mantiene por debajo del
umbral del 5\,\% en $N=14$ y en la práctica totalidad del rango estudiado, con un repunte leve y explicado
(\!$\sim 4,8\,\%$) en $N=18$--$20$ atribuible al límite de expresividad del circuito a profundidad fija
(Figura 6.2).

En conjunto, se considera cumplido también el objetivo general: diseñar, implementar, optimizar y evaluar de
forma sistemática un algoritmo cuántico híbrido para la selección de carteras con restricciones de cardinalidad,
contrastado frente a referencias clásicas industriales.

\section{Limitaciones del Trabajo}
\label{sec:limitaciones}

Pese a los resultados positivos, es necesario reconocer de forma explícita las limitaciones del estudio, tanto por
rigor metodológico como para delimitar correctamente el alcance de las conclusiones anteriores:

\begin{itemize}
    \item \textbf{Simulación sin ruido físico.} Todos los experimentos se han ejecutado mediante simulación exacta
    de vector de estado (\emph{statevector}), sin modelar canales de ruido de hardware real (decoherencia, errores
    de puerta, \emph{crosstalk}). Los resultados de viabilidad y GAP representan, por tanto, una cota superior de
    rendimiento respecto a lo que se observaría en un procesador cuántico físico actual.
    \item \textbf{Contribución real de la regularización Ridge L2.} Como se ha señalado en la Sección
    \ref{sec:cumplimiento_objetivos}, la ablación de $\alpha$ muestra que la penalización Ridge, de forma aislada,
    no mejora el GAP medio e incluso puede empeorarlo ligeramente; su papel en la arquitectura final es más
    limitado de lo planteado inicialmente en el diseño teórico, y requeriría una recalibración de hiperparámetros
    más exhaustiva para confirmar si aporta valor en algún régimen específico.
    \item \textbf{Límite de expresividad a profundidad fija.} El repunte del GAP observado en $N=18$ y $N=20$
    confirma que, para una profundidad de circuito $p$ fija, el ansatz agota su capacidad de representación al
    crecer la dimensión del problema. El estudio de escalado en $p$ (Sección 6.1.3) sugiere que este límite es
    superable aumentando la profundidad, pero dicha vía no ha podido explorarse de forma exhaustiva por el coste
    computacional asociado a la simulación clásica de circuitos más profundos.
    \item \textbf{Coste computacional de la simulación clásica.} El tiempo de ejecución del bucle híbrido
    clásico-cuántico escala de forma notablemente más costosa que los solucionadores clásicos de referencia,
    alcanzando cientos de segundos por instancia para $N$ elevado. Esto limita la escala de universos de activos
    evaluables en este trabajo y aleja, por ahora, la viabilidad de un despliegue en tiempo real.
    \item \textbf{Alcance del modelo financiero.} El estudio se restringe a un modelo estático de media-varianza de
    Markowitz de un único periodo, sobre un universo de activos de renta variable (S\&P 500 / IBEX). No se han
    modelado costes de transacción, rebalanceo dinámico de cartera, ni otras clases de activos, por lo que la
    generalización de los resultados a escenarios de gestión de carteras más complejos queda pendiente de
    validación.
    \item \textbf{Sensibilidad del optimizador clásico.} El bucle de optimización clásica (COBYLA/SLSQP) puede verse
    afectado por mesetas estériles (\emph{barren plateaus}) al escalar el número de cúbits; si bien la inicialización
    TQA mitiga este riesgo en el rango estudiado, no se ha caracterizado su comportamiento para universos de activos
    sustancialmente mayores a $N=20$.
\end{itemize}

\section{Líneas Futuras de Trabajo}
\label{sec:trabajo_futuro}

A partir de las limitaciones anteriores y de los hallazgos del proyecto, se identifican las siguientes líneas de
continuación, ordenadas de menor a mayor alcance:

\begin{enumerate}
    \item \textbf{Validación en hardware cuántico real.} Ejecutar el circuito XY-QAOA sobre plataformas NISQ físicas
    (p.\,ej. IBM Quantum, IonQ) incorporando técnicas de mitigación de errores, para contrastar la degradación de
    viabilidad y GAP frente al comportamiento observado en simulación exacta y cuantificar la brecha
    simulación--hardware real.
    \item \textbf{Recalibración sistemática de la regularización.} Dado que la ablación de $\alpha$ ha revelado que
    el efecto de Ridge L2 es más limitado de lo previsto, explorar una búsqueda de hiperparámetros más amplia
    (por ejemplo, mediante optimización bayesiana) sobre $\alpha$, la rampa temporal $T$ de TQA y el número de
    reinicios, para determinar si existe una región del espacio de hiperparámetros donde la penalización aporte una
    mejora consistente.
    \item \textbf{Estrategias de profundidad adaptativa.} Investigar esquemas de \emph{warm-start} o de crecimiento
    incremental de la profundidad $p$ (por ejemplo, inspirados en ADAPT-QAOA) que permitan compensar el límite de
    expresividad detectado en instancias grandes sin incurrir en el coste de fijar una profundidad elevada desde el
    inicio.
    \item \textbf{Extensión del modelo financiero.} Incorporar rebalanceo dinámico multiperiodo, costes de
    transacción y restricciones sectoriales o de liquidez, acercando el modelo QUBO a escenarios de gestión de
    carteras más realistas y exigentes desde el punto de vista de la industria.
    \item \textbf{Ampliación del universo de activos y regímenes de mercado.} Evaluar el algoritmo sobre clases de
    activos adicionales (renta fija, materias primas, criptoactivos) y bajo distintos regímenes de mercado
    (alta volatilidad, crisis, mercados alcistas), para valorar la robustez de la generalización fuera de muestra
    observada en este trabajo en condiciones más adversas.
    \item \textbf{Exploración de arquitecturas alternativas.} Comparar el XY-Mixer con otros esquemas de
    preservación de restricciones (p.\,ej. \emph{Grover-mixer}, QAOA+ multiángulo) o con enfoques de aprendizaje
    automático cuántico, para situar los resultados de este TFM dentro de un marco comparativo más amplio del
    estado del arte en optimización combinatoria cuántica aplicada a finanzas.
    \item \textbf{Análisis de coste energético a escala.} Extender el análisis de huella de carbono iniciado en el
    Capítulo 5 a escenarios de despliegue en QPU real, comparando el coste energético del enfoque cuántico frente a
    la infraestructura clásica equivalente a medida que crece el tamaño del problema.
\end{enumerate}

En definitiva, este trabajo demuestra que la vía de la preservación estructural de restricciones mediante circuitos
cuánticos especializados constituye una dirección de investigación sólida y con evidencia empírica favorable, si
bien su madurez para un despliegue productivo en el sector financiero depende todavía de avances tanto en el
hardware cuántico disponible como en la sofisticación de los modelos financieros que se le acoplen.